\section{The conductor obstruction}\label{sec:conductor}

We first establish the local algebra needed to distinguish failure of
$S_2$ from nonnormality. Normalization is finite for domains essentially
of finite type over $\C$ \cite[Tags~035B and~035S]{stacks-project}.
If $A$ is such a domain and $B$ is its
normalization in $K=\Frac(A)$, its conductor is
\begin{equation}\label{eq:conductor}
 I=(A:B)=\{a\in A:aB\subseteq A\}.
\end{equation}
It is a nonzero ideal of both $A$ and $B$. Indeed, clearing denominators
of finitely many $A$-module generators of $B$ gives a nonzero element
of $I$. The equality $I=\Ann_A(B/A)$ identifies
$V(I)=\operatorname{Supp}_A(B/A)$ with the locus where normalization
is not an isomorphism. We call this closed set the conductor locus.

\begin{lemma}[A conductor divisor under $S_2$]\label{lem:conductor}
Let $A$ be a local domain essentially of finite type over $\C$,
and assume that $A$ satisfies $S_2$. Let $B$ be its finite normalization
and $I=(A:B)$. If $\mathfrak q\in\Spec B$ satisfies
$I B_{\mathfrak q}\neq B_{\mathfrak q}$, then there is a height-one
prime $\mathfrak h\subseteq\mathfrak q$ of $B$ containing $I$.
\end{lemma}

\begin{proof}
As a $B$-module,
\begin{equation}\label{eq:conductor-hom}
 I\simeq\Hom_A(B,A).
\end{equation}
For completeness, tensoring an $A$-linear map $B\to A$ with $K$
gives a $K$-linear endomorphism of $K$, hence multiplication by some
$u\in K$. Evaluating at $1$ shows $u\in A$, and the requirement
$uB\subseteq A$ gives $u\in I$. Conversely multiplication by every
$u\in I$ defines such a map. The natural $B$-action on the Hom module
corresponds to the ideal structure in~\eqref{eq:conductor-hom}.

We show that $I$ is an $S_2$ module over $B$. Fix
$\mathfrak q\in\Spec B$ and put $\mathfrak p=\mathfrak q\cap A$.
The heights of these primes agree. To check this in the present setting,
write the extension as a localization of the normalization of an affine
domain. The two affine domains have the same fraction field, and the
residue fields at corresponding primes form a finite algebraic extension.
The formula
\[
 \height\mathfrak p=\operatorname{trdeg}_{\C}K
                 -\operatorname{trdeg}_{\C}\kappa(\mathfrak p)
\]
and its counterpart in the normalization therefore give equal heights
\cite[Tags~00OT, 00P0 and~00P1]{stacks-project}.
The localizations involved do not remove any prime below the primes
being considered.

If this common height is at least two, $S_2$ supplies an
$A_{\mathfrak p}$-regular sequence $x,y\in\mathfrak p A_{\mathfrak p}$.
Write $B_{\mathfrak p}$ for localization as an $A$-algebra. Finite
presentation gives
\[
 I_{\mathfrak p}
   =\Hom_{A_{\mathfrak p}}(B_{\mathfrak p},A_{\mathfrak p}).
\]
Multiplication by $x$ is injective on this module. Suppose that
$y\varphi=x\psi$ for two elements of it. For every
$b\in B_{\mathfrak p}$ one has
$y\varphi(b)\in xA_{\mathfrak p}$. Regularity of $y$ modulo $x$
implies $\varphi(b)\in xA_{\mathfrak p}$. Division by $x$ defines
an $A_{\mathfrak p}$-linear map $\eta:B_{\mathfrak p}\to A_{\mathfrak p}$,
with $\varphi=x\eta$. Hence $y$ is a nonzerodivisor modulo $x$ on
$I_{\mathfrak p}$. Localizing further at $\mathfrak q$ preserves this
regular sequence; both elements remain in the maximal ideal there.
In height one, torsion-freeness of the nonzero ideal $I$ gives depth
at least one. This proves the asserted $S_2$ property of $I$.

Choose a prime $\mathfrak h$ minimal over $I$ and contained in
$\mathfrak q$. Its height is positive because $I$ is a nonzero ideal
in a domain. If $\height\mathfrak h\geq2$, normality of $B$ and the preceding
paragraph imply
\[
 \depth B_{\mathfrak h}\geq2,
 \qquad \depth I_{\mathfrak h}\geq2.
\]
Apply the depth lemma \cite[Tag~00LX]{stacks-project} to
\[
 0\longrightarrow I_{\mathfrak h}
  \longrightarrow B_{\mathfrak h}
  \longrightarrow (B/I)_{\mathfrak h}\longrightarrow0.
\]
It gives $\depth(B/I)_{\mathfrak h}\geq1$. Minimality of
$\mathfrak h$ over $I$, however, makes $(B/I)_{\mathfrak h}$ a nonzero
zero-dimensional local module. Such a module has depth zero.
This contradiction proves $\height\mathfrak h=1$.
\end{proof}

\begin{lemma}\label{lem:proper-at-fiber}
Under the notation of~\eqref{eq:conductor}, suppose $A$ is local with
maximal ideal $\mathfrak n$ and that $\Spec B$ has at least two points
above $\mathfrak n$. Then the conductor is proper at every point above
$\mathfrak n$.
\end{lemma}

\begin{proof}
If $I B_{\mathfrak q}=B_{\mathfrak q}$ at a prime above
$\mathfrak n$, some $c\in I$ avoids $\mathfrak q$. Since $c\in A$
and $\mathfrak q\cap A=\mathfrak n$, it is a unit in $A$. Thus
$I=A$, so $B=A$, contrary to the existence of the two points.
\end{proof}

\begin{proposition}[The two-point obstruction]\label{prop:conductor-obstruction}
Let $X$ be an irreducible affine complex variety and let $f:Y\to X$
be a dominant birational morphism, with $Y$ smooth and irreducible. Suppose there is an open subset $G\subseteq Y$
such that
\[
 f(G)\subseteq X_{\mathrm{sm}},\qquad
 \codim_Y(Y\setminus G)\geq2.
\]
If distinct closed points $y_1,y_2\in Y$ have the same image $T\in X$
and both differentials $df_{y_i}$ are injective, then
$\OO_{X,T}$ does not satisfy $S_2$ as a ring and is not
Cohen--Macaulay. In particular $X$ does not satisfy $S_2$.
\end{proposition}

\begin{proof}
Let $\nu:\overline X\to X$ be the finite normalization.
The dominant morphism $f$ lifts to a morphism
\[
 g:Y\longrightarrow\overline X,
 \qquad \nu\circ g=f.
\]
Here $Y$ is normal because it is smooth. The normalization lifting
property is \cite[Tag~035Q(4)]{stacks-project}; in this setting the
property follows directly from integral closure: pullbacks of integral
rational functions are integral over the relevant local rings of $Y$
and hence regular. Since $f$ and $\nu$ are birational, $g$ is birational as well.

Injectivity of $df_{y_i}$ makes the Zariski tangent space of the fiber
$f^{-1}(T)$ zero at $y_i$. The local ring of that fiber therefore has
dimension zero, so $f$ is quasi-finite at $y_i$. The fiber of $g$ through
$y_i$ is contained in that fiber, and $g$ is also quasi-finite there.
Let $Y_0\subseteq Y$ be the open quasi-finite locus of $g$, containing
both $y_i$ \cite[Tag~01TI]{stacks-project}.
The morphism $g|_{Y_0}$ is separated, quasi-finite and
birational with normal target. Since $\overline X$ is affine,
Zariski's Main Theorem factors it as an open immersion
$Y_0\hookrightarrow Z$ followed by a finite morphism $Z\to\overline X$
\cite[Tag~05K0]{stacks-project}; replacing $Z$ by the closure of the
image of the integral scheme $Y_0$, with its reduced structure, we may
take $Z$ integral. Then $Z\to\overline X$ is finite and birational with
normal target, hence an isomorphism \cite[Tag~0AB1]{stacks-project},
so $g|_{Y_0}$ is an open immersion.
In particular,
\[
 q_i=g(y_i)\in\overline X
\]
are distinct points lying over $T$, and an open neighborhood of each
$q_i$ is identified with the corresponding open subset of $Y$.

Assume for contradiction that $A=\OO_{X,T}$ satisfies $S_2$.
Its normalization $B$ is the localization of the affine normalization
above $T$, and the points $q_i$ give two distinct primes of $B$ above
the maximal ideal of $A$. The conductor also commutes with this
localization: use~\eqref{eq:conductor-hom} and the fact that the finite
normalization is a finitely presented module over the Noetherian
coordinate ring. Lemmas~\ref{lem:conductor} and
\ref{lem:proper-at-fiber} give a height-one prime in the conductor
locus whose closure passes through $q_1$.

The same height-one prime defines a divisor on a neighborhood of
$q_1$ in the affine normalization. Indeed, localization from the
affine normalization to $B$ preserves heights for all primes below
$q_1$. The open chart containing $q_1$ also contains that height-one
generic point. Consequently the chart $g(Y_0)$ contains a conductor
divisor.

On the other hand, $f(G)$ is contained in the smooth locus of $X$
by hypothesis. Normalization is an isomorphism over that locus, so the conductor locus is disjoint from
$g(Y_0\cap G)$. Under the open-chart identification,
this locus is therefore contained in
\[
 Y_0\setminus G.
\]
This set has codimension at least two by hypothesis;
it cannot contain the divisor just obtained. This contradiction
proves that $A$ is not an $S_2$ ring. A Cohen--Macaulay local ring
has Cohen--Macaulay localizations
\cite[Tag~00NB]{stacks-project} and hence is $S_2$, so $A$ is not
Cohen--Macaulay. If $X$ were $S_2$, its localization $A$ would
be $S_2$ as well, giving the final assertion.
\end{proof}

\begin{remark}\label{rem:fiber-scope}
Only the two specified immersive decompositions are used. Other
decompositions, additional normalization points, and positive-dimensional
fibers elsewhere cause no difficulty: Zariski's Main Theorem is applied
only on the open quasi-finite locus. The argument also makes no claim
that $X$ is smooth in codimension one.
\end{remark}
